\documentclass[border=5pt]{standalone}
% Style commun aux figures de structure de MPMbox
% ------------------------------------------------
\usepackage[utf8]{inputenc}
\usepackage[T1]{fontenc}
\usepackage[french]{babel}
\usepackage{lmodern}
\usepackage{tikz}
\usetikzlibrary{positioning,arrows.meta,calc,fit,backgrounds,shapes.geometric}
% babel(french) rend « : » actif, ce qui casse les coordonnées polaires (90:4cm)
\usetikzlibrary{babel}

% pas de césure dans les boîtes étroites
\hyphenpenalty=10000
\exhyphenpenalty=10000
\tolerance=9999
\emergencystretch=3em

\definecolor{mpmBase}{HTML}{1F3A5F}    % classes de base
\definecolor{mpmBaseBg}{HTML}{E8EEF6}
\definecolor{mpmDeriv}{HTML}{4A7BA7}   % variantes
\definecolor{mpmDerivBg}{HTML}{F5F8FB}
\definecolor{mpmData}{HTML}{7A5230}    % données simulées
\definecolor{mpmDataBg}{HTML}{FBF3E8}
\definecolor{mpmAccent}{HTML}{A6432F}  % mise en avant (MPMxDEM)
\definecolor{mpmAccentBg}{HTML}{FBEEEB}
\definecolor{mpmGrey}{HTML}{5A6472}

\tikzset{
  base/.style={
    draw=mpmBase, line width=0.9pt, fill=mpmBaseBg, rounded corners=3pt,
    text width=#1, align=center, inner sep=6pt, font=\sffamily\small},
  base/.default=5.2cm,
  deriv/.style={
    draw=mpmDeriv, line width=0.6pt, fill=mpmDerivBg, rounded corners=2pt,
    text width=#1, align=left, inner sep=5pt, font=\sffamily\scriptsize},
  deriv/.default=7.4cm,
  data/.style={
    draw=mpmData, line width=0.7pt, fill=mpmDataBg, rounded corners=3pt,
    text width=#1, align=center, inner sep=5pt, font=\sffamily\scriptsize},
  data/.default=4.2cm,
  accent/.style={
    draw=mpmAccent, line width=0.8pt, fill=mpmAccentBg, rounded corners=2pt,
    text width=#1, align=left, inner sep=5pt, font=\sffamily\scriptsize},
  accent/.default=7.4cm,
  spine/.style={draw=mpmDeriv, line width=0.6pt},
  link/.style={draw=mpmGrey, line width=0.6pt, -{Stealth[length=4pt]}},
  dashlink/.style={draw=mpmGrey, line width=0.6pt, densely dashed,
                   -{Stealth[length=4pt]}, font=\sffamily\tiny},
  mpmcap/.style={font=\sffamily\scriptsize\itshape, text=mpmGrey},
}

% Titre d'une classe de base : \classname{Nom}{rôle}
\newcommand{\classname}[2]{{\bfseries #1}\\[2pt]{\scriptsize\itshape #2}}
% Ligne d'une variante : \variant{Nom}{description}
\newcommand{\variant}[2]{{\bfseries\ttfamily #1}\quad #2}

\begin{document}
\begin{tikzpicture}

% ---- le chef d'orchestre
\node[base=4.4cm] (BOX) at (0,0)
  {\classname{MPMbox}{lit les données,\\ boucle en temps,\\ écrit les résultats}};

% ---- ce qui est simulé
\node[data=4.8cm] (MP)   at (-6.2, 2.1) {{\bfseries MaterialPoint}\\ points matériels porteurs de la matière\\ (masse, vitesse, contrainte)};
\node[data=4.4cm] (GRID) at (-6.2, 0.0) {{\bfseries grid / node / element}\\ grille fixe de calcul,\\ remise à zéro à chaque pas};
\node[data=4.4cm] (OBS)  at (-6.2,-2.1) {{\bfseries Obstacle}\\ corps rigides qui bornent\\ ou pilotent l'écoulement};

% ---- les ingrédients interchangeables
\node[base=4.0cm] (CM)  at (6.0, 2.55) {{\bfseries ConstitutiveModel}\\[1pt]{\scriptsize\itshape comportement du matériau\\ un par point matériel}};
\node[base=4.0cm] (SF)  at (6.0, 0.85) {{\bfseries ShapeFunction}\\[1pt]{\scriptsize\itshape transfert points $\leftrightarrow$ grille}};
\node[base=4.0cm] (OS)  at (6.0,-0.85) {{\bfseries OneStep}\\[1pt]{\scriptsize\itshape schéma d'intégration}};
\node[base=4.0cm] (BFL) at (6.0,-2.55) {{\bfseries BoundaryForceLaw}\\[1pt]{\scriptsize\itshape loi de contact\\ une par obstacle}};

\node[base=4.0cm] (CMD) at (10.8, 2.55) {{\bfseries Command}\\[1pt]{\scriptsize\itshape mise en données}};
\node[base=4.0cm] (SCH) at (10.8, 0.85) {{\bfseries Scheduler}\\[1pt]{\scriptsize\itshape scénario au cours du temps}};
\node[base=4.0cm] (SPY) at (10.8,-0.85) {{\bfseries Spy}\\[1pt]{\scriptsize\itshape mesures pendant le calcul}};
% ---- regroupements
\begin{scope}[on background layer]
  \node[draw=mpmData, densely dashed, rounded corners=4pt, inner sep=9pt,
        fit=(MP)(GRID)(OBS)] (GDATA) {};
  \node[draw=mpmBase, densely dashed, rounded corners=4pt, inner sep=9pt,
        fit=(CM)(BFL)(CMD)(SPY)] (GCHOIX) {};
\end{scope}

\node[font=\sffamily\small\bfseries, text=mpmData, anchor=south]
  at ($(GDATA.north)+(0,3pt)$) {ce qui est simulé};
\node[font=\sffamily\small\bfseries, text=mpmBase, anchor=south]
  at ($(GCHOIX.north)+(0,3pt)$) {ce que l'utilisateur choisit, par un mot-clé du fichier d'entrée};

\draw[link, line width=1pt] (BOX.west)  -- (GDATA.east);
\draw[link, line width=1pt] (BOX.east)  -- (GCHOIX.west);

\node[mpmcap, anchor=north, text width=13cm, align=center] at ($(GCHOIX.south)+(-3.0cm,-0.35cm)$)
  {le cœur du code ne connaît que les familles de gauche et de droite,
   jamais la variante réellement choisie};

\end{tikzpicture}
\end{document}
